\documentclass[12pt,a4paper]{article}

\usepackage[utf8]{inputenc}
\usepackage[T1]{fontenc}
\usepackage[brazil]{babel}
\usepackage{setspace}
\usepackage{geometry}
\geometry{margin=2.5cm}
\setstretch{1.25}

\title{Beneath the Surface: Investigating Anti-Competitive Signals in Brazilian Public Procurement}

\author{
  Sergio Galletta\\
  \small ETH Z\"urich
  \and
  Darcio Genicolo-Martins\\
  \small Insper -- Instituto de Ensino e Pesquisa
  \and
  Tommaso Giommoni\\
  \small University of Amsterdam
}

\date{}

\begin{document}

\maketitle

O estudo investiga indícios de comportamentos anticompetitivos em licitações públicas do estado de São Paulo utilizando uma estratégia empírica inspirada em Kawai et al.\ (2023), mas adaptada a um ambiente substancialmente mais complexo, heterogêneo e institucionalmente rico do que o cenário original analisado pelos autores. A motivação central é simples e poderosa: quando empresas disputam um contrato público e a diferença entre o lance vencedor e o imediatamente inferior é ínfima — por vezes, meros centésimos ou milésimos — o resultado deveria refletir exclusivamente a competição de preços baseada em custos marginais, eficiência operacional e choques idiossincráticos de pequena magnitude. Essa zona de ``near tie'' funciona como uma espécie de laboratório natural de competição quase perfeita: atores semelhantes, disputando um objeto bem definido, sob regras de mercado explícitas e observáveis. Em teoria, nenhuma empresa deveria ser particularmente mais propensa a vencer quando a margem que a separa de sua rival é tão pequena que sua importância econômica se aproxima de zero. Se, ao cruzar esse limite, observamos diferenças persistentes e sistemáticas no perfil das empresas, no seu histórico de vitórias, em sua forma de participar da licitação ou em seus padrões de envio de lances, então essas diferenças dificilmente podem ser atribuídas à aleatoriedade competitiva. Elas revelam, ao contrário, que algum mecanismo adicional está operando, seja ele colusão explícita entre participantes, coordenação tácita sustentada por repetição e monitoramento, ou favorecimento institucional originado no lado da administração pública.

Essa capacidade de extrair sinais estruturais a partir de margens mínimas torna o método particularmente adequado para ambientes onde a transparência formal é elevada, mas a competição real é menor do que aparenta. Em plataformas eletrônicas como a Bolsa Eletrônica de Compras (BEC), o desenho institucional comunica uma mensagem de neutralidade: regras claras, ampla publicidade, múltiplos participantes, registros detalhados e rastreáveis. Contudo, isso não impede que, por baixo dessa superfície de transparência, se desenvolvam arranjos informais de coordenação entre firmas ou relações privilegiadas entre alguns fornecedores e determinados compradores. Nesses ambientes, práticas colusivas tendem a se sofisticar: em vez de sinais grosseiros, como preços extremamente acima do mercado, observam-se sutilezas comportamentais, como pequenas alternâncias estratégicas de vitórias, padrões de entrada cuidadosamente coordenados ou uso cirúrgico da informação sobre o comportamento dos rivais. O foco em near ties permite justamente isolar esse tipo de sutileza, ao concentrar a análise em situações em que a competição deveria ser mais ``pura'' e o espaço para manipulação de resultado, se existir, será especialmente revelador.

A base de dados utilizada reúne mais de 3,8 milhões de lances enviados por aproximadamente 19 mil firmas ao longo de uma década inteira de compras governamentais no Estado de São Paulo, cobrindo o período de 2009 a 2019. Cada observação registra não apenas o valor ofertado e o resultado do leilão, mas também informações sobre o item, o lote, a unidade compradora, o tipo de procedimento, o valor de referência estabelecido pela administração, o momento do envio de cada lance e a sequência completa de movimentos realizados pelos participantes. A BEC constitui, assim, não apenas um canal de contratação pública, mas também uma fonte riquíssima de dados administrativos que documentam, em alta resolução temporal, a interação entre Estado e fornecedores privados. Essa riqueza de informação nos permite construir, para cada firma, uma trajetória detalhada de participação: podemos observar quais mercados ela frequenta, com que intensidade, com quais rivais ela se encontra com maior frequência, em que medida suas vitórias são concentradas em alguns segmentos específicos ou dispersas ao longo de múltiplos mercados.

Essa granularidade também viabiliza a construção de medidas de interação entre firmas ao longo do tempo, aproximando a análise de enfoque de redes de relacionamento. É possível, por exemplo, identificar subconjuntos de empresas que aparecem juntas repetidamente em licitações similares, o que abre espaço para interpretar os padrões de near ties não apenas sob a ótica de firmas individuais, mas também sob a perspectiva de grupos ou ``anéis'' de competidores recorrentes. Em um cenário de competição efetiva, a presença reiterada das mesmas firmas num mesmo mercado pode simplesmente refletir especialização e vantagens comparativas. Em um cenário de conluio, por outro lado, a recorrência de combinações específicas de participantes em near ties pode ser indício de um pacto estável de divisão de mercado, onde os papéis de vencedor e perdedor são rotineiramente alternados.

A estratégia de identificação se baseia na comparação entre empresas que ganham e perdem por margens extremamente estreitas, explorando uma descontinuidade no \textit{cutoff} zero da diferença relativa entre o primeiro e o segundo melhor lance. Essa diferença funciona como a variável contínua central do desenho de \textit{Regression Discontinuity} (RDD). Quando a diferença relativa se aproxima de zero, a teoria competitiva sugere que o resultado é determinado por ruído, isto é, por fatores que não guardam relação sistemática com atributos estruturais das firmas. Ao restringir a amostra a esse subconjunto de disputas, construímos uma situação em que, sob competição, a probabilidade de vitória para cada participante deveria ser praticamente proporcional apenas à sua presença ex ante, e não a fatores como incumbência, conexões políticas ou conhecimento prévio dos lances rivais. Em outras palavras, a linha que separa quem ganha de quem perde em near ties deveria ser essencialmente arbitrária.

Esse argumento de arbitrariedade é o núcleo da interpretação causal do desenho RDD. Quando, empiricamente, detectamos que a passagem de perdas mínimas para vitórias mínimas está associada a saltos discretos em variáveis que refletem o estágio em que a firma se encontra no relacionamento com o Estado (como incumbência) ou na forma como ela constrói sua estratégia de lances (como ser o último a ofertar), somos levados a concluir que tais variáveis estão diretamente vinculadas a mecanismos não competitivos que operam no ponto de decisão. Vale ressaltar que essa lógica não exige que observemos diretamente comunicação entre firmas ou agentes públicos: basta que padrões de comportamento sejam suficientemente sistemáticos e descontinuamente alinhados com a linha divisória entre vitória e derrota para que possamos rejeitar a hipótese de pura competição. Essa é uma virtude central da abordagem: ela transforma informações rotineiramente geradas pelos sistemas de compras — e muitas vezes pouco exploradas — em evidências estruturais sobre o funcionamento do mercado.

A análise empírica concentra-se em dois indicadores-chave: (i) incumbência, definida como ter vencido a licitação anterior naquele mesmo mercado, e (ii) a probabilidade de enviar o último lance. A definição de mercado que utilizamos segue uma lógica que combina o tipo de bem ou serviço com o conjunto de firmas que interagem em licitações anteriores, formando grupos de competição relativamente estáveis ao longo do tempo. Assim, ser incumbente significa não apenas ter ganho um contrato passado, mas ter ganhado em um contexto em que as mesmas firmas voltarão a se encontrar, criando condições para monitoramento mútuo e potencial disciplina de cartel. Em mercados repetidos, a incumbência pode ser usada como mecanismo de recompensa por cooperação passada ou de punição por desvios: firmas que ``respeitam'' acordos implícitos de não agressão em determinados leilões podem ser favorecidas em rodadas futuras, enquanto firmas que tentam romper a disciplina, ofertando preços abaixo do combinado, podem ser excluídas de ciclos de vitórias ou alvo de retaliação.

Já o indicador de último lance captura a dimensão temporal da estratégia da firma dentro do leilão específico. Em muitos formatos de licitação, especialmente aqueles com rodadas abertas ou com períodos finais de ofertas, o momento exato em que o lance é enviado pode transmitir informações sobre o comportamento do licitante: lances enviados muito cedo podem indicar tentativa de ancorar o processo, enquanto lances enviados muito tarde podem sinalizar que a firma está reagindo a um fluxo de informações que não é igualmente acessível a todos os participantes. Se uma firma com frequência posiciona seu lance final de forma a vencer por uma margem mínima, isso sugere que ela tem alguma clareza sobre a posição dos rivais, seja por observar diretamente seus comportamentos, seja por receber sinais indiretos. Quando essa característica — ser o último a ofertar — aparece de forma desproporcional entre os vencedores por pouco, e quase não se manifesta entre os perdedores por pouco, a hipótese de \textit{bid leakage} ou de coordenação torna-se difícil de ignorar.

Os resultados empíricos mostram descontinuidades claras e estatisticamente significativas em ambos os indicadores no \textit{cutoff} da margem de vitória. Para incumbência, encontramos um aumento de aproximadamente 1,8 ponto percentual na probabilidade de o vencedor por margem mínima ser incumbente, mesmo após controlar para uma ampla gama de efeitos fixos e variáveis de mercado. Esse resultado é particularmente impressionante quando consideramos o tamanho da amostra e o caráter extremamente local da variação explorada: estamos falando de diferenças que aparecem em um intervalo minúsculo da distribuição de lances, o que torna improvável que estejamos capturando apenas correlações espúrias ou tendências de longo prazo. Em um mundo puramente competitivo, a probabilidade de incumbência entre vencedores e perdedores em near ties deveria ser aproximadamente a mesma. O fato de observamos um salto persistente nessa probabilidade indica que incumbentes não apenas se beneficiam de experiência, mas parecem estar inseridos em uma dinâmica de favoritismo estrutural ou coordenação com outras firmas.

Para o indicador do último lance, observamos um salto ainda maior: aproximadamente 6,4 pontos percentuais na probabilidade de o vencedor por pouco ter enviado a última oferta. Esse resultado sugere um padrão recorrente em que a firma que cruza a linha entre derrota e vitória em near ties é, com frequência desproporcional, aquela que ``encerra'' o processo competitivo com um ajuste fino em sua oferta. Essa evidência é consistente com dois tipos principais de mecanismos: (i) arranjos colusivos em que firmas coordenam quem enviará o último lance e em que faixa de valores, de modo a garantir que um membro ``designado'' do grupo vença, ou (ii) vazamento de informações por agentes internos que sinalizam a uma firma específica quando e como ajustar seu lance para garantir a vitória por margem pequena. Em ambos os casos, o que nos interessa é que tais comportamentos não são compatíveis com a hipótese de competição simétrica baseada apenas em custos.

A robustez desses resultados foi extensivamente testada por meio de variações de \textit{bandwidth}, mudanças nos \textit{kernels} utilizados, exclusão de \textit{outliers}, uso de janelas simétricas e assimétricas, inclusão de controles adicionais sobre mercado e tempo, e testes de continuidade sobre a densidade da variável contínua. Em particular, aplicamos rotinas padrão da literatura de RDD para verificar se existe manipulação local da variável running, isto é, se firmas parecem ``empurrar'' seus lances de forma a se concentrar em torno de determinados pontos da distribuição. Não encontramos evidência robusta de manipulação sistemática da margem de vitória em si, mas sim de que, dado um conjunto de near ties, a composição de quem vence e quem perde muda de forma abrupta em termos de incumbência e padrão de timing. Essa combinação — ausência de manipulação da variável running e presença de descontinuidade em variáveis de resultado — reforça a interpretação causal das estimativas.

Além das análises de linha de base, exploramos heterogeneidades em diversos recortes, como o número de participantes, o tipo de item, a frequência de licitações em determinado mercado e a presença de histórico conhecido de irregularidades formais em certos órgãos ou segmentos. De maneira geral, encontramos que os sinais de descontinuidade são mais fortes em mercados com poucos participantes recorrentes, o que é consistente com a teoria de que cartéis são mais fáceis de sustentar quando há menor rotatividade de rivais. Também observamos que alguns grupos de itens, especialmente aqueles caracterizados por elevada especificidade técnica ou por histórico de uso intensivo por determinados órgãos, exibem padrões mais intensos de incumbência em near ties. Esses resultados de heterogeneidade sugerem que os mecanismos não competitivos detectados não são uniformes, mas se concentram em nichos específicos do sistema de compras, o que é informativo para o desenho de políticas de fiscalização mais direcionadas.

A interpretação econômica dos achados se alinha firmemente com modelos teóricos de coordenação e conluio em mercados repetidos. Cartéis costumam operar em ambientes onde as empresas têm oportunidades frequentes de interagir, o que cria condições favoráveis para pactos de não agressão, alternância de vencedores e disciplina interna por meio de estratégias de gatilho. No caso da BEC, que reúne milhares de licitações similares ao longo do tempo, as oportunidades para monitoramento mútuo e punição de desvios são amplas. Além disso, a padronização dos itens facilita a criação de expectativas estáveis sobre custos, margens e comportamento dos rivais, reduzindo incertezas que poderiam comprometer a sustentação de acordos colusivos. O salto na incumbência é compatível com firmas que, em vez de competir agressivamente, seguem uma regra tácita de revezamento ou de permanência prolongada em posição dominante, na qual vitórias anteriores condicionam vitórias futuras. Essa configuração é congruente com histórias de ``campeões crônicos'', firmes que concentram contratos em determinados segmentos, mesmo em contextos de aparente multiprovedoria.

Já o padrão do último lance se aproxima de situações em que a firma vencedora tem acesso privilegiado a informações, permitindo-lhe ajustar seu lance final de maneira precisa para vencer por uma margem mínima — um comportamento característico de arranjos sofisticados de conluio ou de vazamentos de informação do lado da administração pública. Em muitos esquemas de \textit{bid leakage} documentados na prática, o agente público que conduz a licitação informa a firma favorecida sobre o melhor lance dos concorrentes, permitindo que ela submeta uma proposta ligeiramente inferior e, assim, vença. Em um ambiente eletrônico, isso pode se manifestar em padrões de última oferta perfeitamente calibrados. Encontrar esse tipo de padrão de forma recorrente em uma base extensa como a nossa reforça a hipótese de que há, em alguns mercados, uma estrutura de relacionamento entre fornecedores e compradores que ultrapassa o modelo padrão de competição impessoal em plataformas digitais.

Além disso, as evidências encontradas têm implicações práticas e de política pública consideráveis. A metodologia apresentada no estudo pode ser facilmente incorporada em sistemas de auditoria preventiva de órgãos de controle, uma vez que os cálculos necessários são simples, replicáveis e escaláveis. Um algoritmo pode monitorar diariamente a ocorrência de near ties, calcular as taxas de incumbência e de último lance em torno do \textit{cutoff} e identificar padrões anômalos em mercados específicos, sinalizando quais deles merecem investigação mais aprofundada. Essa prática permitiria que controladorias, tribunais de contas e ministérios públicos atuassem de forma mais proativa, focando recursos em segmentos do mercado onde o risco de conluio ou favorecimento é empiricamente maior. Em vez de depender exclusivamente de denúncias pontuais, muitas vezes tardias e difíceis de comprovar, as instituições de controle passariam a contar com uma triagem baseada em dados, capaz de apontar onde os indícios comportamentais mais fortes se concentram.

A adoção dessa abordagem também incentiva uma mudança na forma como os próprios sistemas de compras são desenhados. Ao revelar que determinados formatos de licitação ou regras procedimentais podem ser mais vulneráveis a práticas anticompetitivas, o método fornece insumos concretos para a reforma de regulamentos, manuais de boas práticas e parâmetros de monitoramento. Por exemplo, se observarmos que mercados com maior flexibilidade temporal para envio de lances exibem mais sinais de uso estratégico do último lance em near ties, isso pode motivar repensar janelas de tempo, mecanismos de encerramento e protocolos de anonimização de lances. De maneira análoga, se segmentos com forte concentração de incumbência em near ties são identificados como de alto risco, o poder público pode adotar medidas específicas, como estímulos à entrada de novos ofertantes, revisão de especificações técnicas, uso de leilões reversos mais agressivos ou fiscalização especial de contratos.

\bigskip
\textbf{Contribuição à literatura.} Este estudo representa uma contribuição substancial à literatura sobre detecção de conluio, economia de compras públicas e métodos semiparamétricos de identificação causal. Ao adaptar e expandir a metodologia de \textit{near-tie Regression Discontinuity} em um ambiente institucional vasto, diversificado e altamente digitalizado, nosso trabalho demonstra que efeitos colusivos podem ser detectados mesmo quando as firmas se empenham em mascarar seu comportamento cooperativo por meio de múltiplos participantes aparentes, dispersão moderada de preços e uso intensivo de plataformas eletrônicas. Diferentemente de abordagens que dependem de anomalias extremas de preços, de casos-limite ou de modelos estruturais complexos de custos, nossa estratégia utiliza uma lógica simples, elegante e baseada em uma variação local muito clara, facilmente comunicável a públicos não especialistas. Essa propriedade torna o método robusto e intuitivamente compreensível para pesquisadores, formuladores de políticas e auditores, reduzindo a distância entre ferramentas sofisticadas de econometria aplicada e o dia a dia das instituições de controle.

Adicionalmente, o estudo amplia a fronteira da literatura ao aplicar esse método a um dos maiores e mais diversificados mercados públicos da América Latina, oferecendo insights que transcendem o caso brasileiro e podem ser generalizados, com as devidas adaptações, a outros países que adotam plataformas eletrônicas de compras. Os resultados encontrados sugerem que comportamentos colusivos podem emergir ou se adaptar mesmo em sistemas altamente informatizados, indicando que a mera transição para plataformas digitais não é suficiente para eliminar práticas anticompetitivas. Isso abre novas frentes de pesquisa sobre como a digitalização afeta a natureza e a detecção do conluio, e motiva estudos futuros que combinem métodos de RDD, análise de redes, \textit{machine learning} e modelos de comportamento estratégico para identificar padrões ainda mais complexos de manipulação e coordenação entre firmas.

Ao demonstrar a aplicabilidade e a força do \textit{near-tie RDD} em um mercado real de grande escala, este artigo contribui tanto para o avanço teórico quanto para o desenvolvimento de ferramentas práticas de combate ao conluio. Do ponto de vista teórico, reforçamos a intuição de que pequenas margens de vitória podem ser usadas como ``janelas'' para observar mecanismos ocultos do jogo competitivo, sobretudo em contextos em que a interação repetida e a observabilidade parcial dos lances permitem sustentação de acordos colusivos sofisticados. Do ponto de vista aplicado, mostramos que esses mesmos mecanismos podem ser transformados em indicadores operacionais simples, que podem ser embutidos em sistemas de TI governamentais e utilizados como parte de uma estratégia contínua de monitoramento de risco. Ao integrar esses dois planos — teoria e prática —, o trabalho contribui para uma agenda de pesquisa que busca não apenas entender como o conluio ocorre, mas também fornecer instrumentos concretos para enfrentá-lo em tempo real.

\end{document}
